\subsection*{Integral Domains}
\label{subsec:integral-domains}

\begin{definition}[Zero Divisor]
\label{def:zero-divisor}
Let $R$ be a ring.
A nonzero element $a\in R$ is a \textbf{zero divisor} if there exists a
nonzero element $b\in R$ such that $ab=0$.
\end{definition}

\begin{remark*}[Standard quantified statement]
\[
\operatorname{ZeroDivisor}(a,R)
\Longleftrightarrow
a\in R
\land a\neq 0
\land
\exists b\in R,\; b\neq 0 \land ab=0.
\]
\end{remark*}

\begin{remark*}[Definition predicate reading]
$a$ is a zero divisor exactly when a nonzero multiplication by $a$ can produce
zero.
\end{remark*}

\begin{remark*}[Negated quantified statement]
\[
\neg\operatorname{ZeroDivisor}(a,R)
\Longleftrightarrow
a=0
\lor
\forall b\in R,\; ab=0 \Rightarrow b=0.
\]
\end{remark*}

\begin{remark*}[Interpretation]
Zero divisors identify where multiplication loses cancellation behavior.
\end{remark*}

\begin{dependencies}
\begin{itemize}
  \item Ring.
  \item Multiplicative zero.
\end{itemize}
\end{dependencies}

\begin{tcolorbox}[title={Definition --- Integral Domain}]
\begin{definition}[Integral Domain]
\label{def:abstract-integral-domain}
An \textbf{integral domain} is a commutative ring with unity in which there are
no zero divisors.
\end{definition}
\end{tcolorbox}

\begin{remark*}[Standard quantified statement]
\[
\operatorname{IntegralDomain}(R,+,\cdot,0,1)
\Longleftrightarrow
\operatorname{CommutativeRing}(R,+,\cdot)
\land
\operatorname{RingWithUnity}(R,+,\cdot,1)
\land
\forall a,b\in R,\; ab=0 \Rightarrow a=0 \lor b=0.
\]
\end{remark*}

\begin{remark*}[Definition predicate reading]
An integral domain is a commutative ring with $1$ where a product is zero only
because at least one factor is zero.
\end{remark*}

\begin{remark*}[Negated quantified statement]
\[
\neg\operatorname{IntegralDomain}(R,+,\cdot,0,1)
\Longleftrightarrow
\neg\operatorname{CommutativeRing}(R,+,\cdot)
\lor
\neg\operatorname{RingWithUnity}(R,+,\cdot,1)
\lor
\exists a,b\in R,\; a\neq0\land b\neq0\land ab=0.
\]
\end{remark*}

\begin{remark*}[Interpretation]
Integral domains are the commutative rings where multiplication has no nonzero
zero-divisor obstruction.
\end{remark*}

\begin{dependencies}
\begin{itemize}
  \item Commutative ring.
  \item Ring with unity.
  \item Zero divisor.
\end{itemize}
\end{dependencies}

\begin{proposition}[Cancellation in Integral Domains]
\label{prop:abstract-domain-cancellation}
Let $R$ be an integral domain.  If $a\neq 0$ and $ab=ac$, then $b=c$.
\hyperref[prf:domain-cancellation]{Go to proof.}
\end{proposition}

\begin{remark*}[Standard quantified statement]
\[
\operatorname{IntegralDomain}(R)
\Longrightarrow
\forall a,b,c\in R,\; a\neq0 \land ab=ac \Rightarrow b=c.
\]
\end{remark*}

\begin{remark*}[Theorem predicate reading]
Nonzero factors can be cancelled in an integral domain.
\end{remark*}

\begin{remark*}[Negated quantified statement]
\[
\operatorname{IntegralDomain}(R)
\land
\exists a,b,c\in R,\; a\neq0 \land ab=ac \land b\neq c
\Longrightarrow \bot.
\]
\end{remark*}

\begin{remark*}[Interpretation]
Cancellation is recovered in integral domains because nonzero elements are not
zero divisors.
\end{remark*}

\begin{dependencies}
\begin{itemize}
  \item Integral domain.
  \item Zero divisor.
\end{itemize}
\end{dependencies}

\begin{remark*}[Proof TODO]
Regenerate this proof as a one-proof-per-file artifact.
\end{remark*}
